% sup2_04_S8a: Supplement part 2, step S8, first half: the Ward-Takahashi identities on a torus and their infinite-lattice limit.
% Input from sup2_00_intro.tex's section (through INPUTS_part2.tex).

\subsection{S8: the Ward--Takahashi identities, and the limit from periodic to free boundary conditions}\label{sup2:ssec:S8}

\cite[pp.~430--431]{Ba3} prints ``the basic identity'' \cite[(2.23)]{Ba3}, its first-order differential \cite[(2.24)]{Ba3} and the examples \cite[(2.25)--(2.28)]{Ba3}, and continues: ``Taking other functions $F$, or differentiating (2.24) to higher order in $A$, we can get all necessary Ward--Takahashi identities. They hold for free boundary conditions also by taking a limit of the identities with periodic boundary conditions.'' (``Free boundary conditions'' is the infinite lattice, as in \cite[p.~584]{Ba4}.) The identities are consumed on \cite[pp.~441--442]{Ba3}, on $\xi\Z^d$ with $C^\xi=(-\Delta^\xi+1)^{-1}$, after ``we resum over proper orderings and $j$-indices, we replace the propagators \dots\ by $C^\xi$'': the square bracket of \cite[(3.29)]{Ba3} is \cite[(2.26)]{Ba3} with $A$ concentrated at one point $y$ and $\lambda(y')=\sum_{\mu'}g'(y)A'_{\mu'}(y)(y'_{\mu'}-y_{\mu'})$, a linear function whose gradient is the constant $g'(y)A'(y)$; the same zeroth-moment identity at the next order is ``apply a proper Ward--Takahashi identity'' in \cite[(3.31)--(3.32)]{Ba3}. So the consumed form is: the identities on the infinite lattice, for full propagator lines, with a constant gauge gradient, for sums over all connected graphs of a fixed order (the vertex positions and the internal photon lines summed). This section proves exactly that.

\paragraph{Setting.} $X$ is the infinite lattice $\eta\Z^d$ or a torus $T_\ell\defeq\eta\Z^d/\ell\Z^d$, $\ell\in\eta\N$. Fields $\phi\colon X\to\R^N$; $q$ real antisymmetric; $U(a)\defeq\exp(e_k\eta qa)$. $B$ is a fixed real bond field (periodic on $T_\ell$), $A$ a bond field, $\langle\phi,-\Delta_{B+A}\phi\rangle\defeq\sum_b\eta^{d-2}|U(B_b+A_b)\phi(b_+)-\phi(b_-)|^2$, and $H(A)\defeq-\Delta_{B+A}+M^2$ with $M^2>0$. For $\lambda\colon X\to\R$ put $(d\lambda)(b)\defeq(\lambda(b_+)-\lambda(b_-))/\eta$ and $(T_\lambda\phi)(x)\defeq\exp(e_kq\lambda(x))\phi(x)$. A polynomial $F$ in finitely many $\phi(x)$ is gauge invariant if $F(T_\lambda\phi)=F(\phi)$ for all $\lambda$; polynomials in $\phi(x)\cdot P(q)\phi(x)$, $P$ a polynomial, are examples, and they include the Wick-ordered $:\!|\phi(x)|^4\!:$, $:\!|\phi(x)|^2\!:$ of \cite[(2.24)--(2.28)]{Ba3} (the Wick matrix is a function of $q$, since every kernel below is built from the $U$) and the mass term $-\frac12\delta m^2:\!|\phi|^2\!:$ ($\delta m^2$ is an even polynomial in $q$, \cite[(1.23)]{Ba3} and \cref{sup2:lem:S7exc}).

\begin{lemma}[exponential decay of the kernels]\label{sup2:lem:S8CT}
Let $H=-\Delta_B+M^2$ on $X$ with any real bond field $B$, and let $C\defeq H^{-1}$ have kernel $C(u,u')$ with respect to the weight $\eta^d$. There are $\theta,c>0$ depending on $d$, $\eta$ and $M$ only (not on $B$, $X$ or $\ell$) such that $|C(u,u')|\le c\,e^{-\theta\dist_X(u,u')}$, $\dist_X$ the sup-distance on $X$.
\end{lemma}

\begin{proof}
$H\ge M^2$ since $-\Delta_B\ge0$. Fix $u'$ and let $\rho(x)\defeq\dist_X(x,u')$, which is $1$-Lipschitz for the sup-distance. The conjugate $H_\theta\defeq e^{\theta\rho}He^{-\theta\rho}$ differs from $H$ only in the hopping terms, each multiplied by $e^{\pm\theta(\rho(b_+)-\rho(b_-))}\in[e^{-\theta\eta},e^{\theta\eta}]$; hence $\norm{H_\theta-H}\le2d\eta^{-2}(e^{\theta\eta}-1)\le M^2/2$ for $\theta\le\theta_0(d,\eta,M)$, and $H_\theta$ is invertible with $\norm{H_\theta^{-1}}\le2/M^2$. Then $C(u,u')=e^{-\theta\rho(u)}(H_\theta^{-1})(u,u')e^{\theta\rho(u')}$ with $\rho(u')=0$, and a kernel entry is bounded by $\eta^{-d}$ times the operator norm.
\end{proof}

\begin{lemma}[gauge covariance]\label{sup2:lem:S8cov}
$H(A+d\lambda)=T_\lambda^{-1}H(A)T_\lambda$ on $X$, for every $\lambda$ (finitely supported if $X=\eta\Z^d$; arbitrary on $T_\ell$).
\end{lemma}

\begin{proof}
With $\psi=T_\lambda\phi$: $U(B_b+A_b+d\lambda(b))(T_\lambda^{-1}\psi)(b_+)-(T_\lambda^{-1}\psi)(b_-)=e^{-e_kq\lambda(b_-)}[U(B_b+A_b)\psi(b_+)-\psi(b_-)]$, because $U(\eta\,d\lambda(b))=\exp(e_kq(\lambda(b_+)-\lambda(b_-)))$ and all these matrices commute. The prefactor is orthogonal, so the quadratic forms agree, and $M^2$ commutes with $T_\lambda$.
\end{proof}

\begin{lemma}[the identities on a torus: exact, every order, connected form]\label{sup2:lem:S8torus}
On $T_\ell$ let $F_1,\dots,F_p$ be gauge-invariant polynomials and, for a multi-index $\beta$, $Z_\beta(A)\defeq\int d\phi\,e^{-\frac12\langle\phi,H(A)\phi\rangle}F^\beta(\phi)$, $F^\beta\defeq\prod_rF_r^{\beta_r}$; let $Z(A,\tau)\defeq\sum_\beta\tau^\beta Z_\beta(A)/\beta!$ as a formal power series in $\tau$. Then $Z_\beta(A+d\lambda)=Z_\beta(A)$ for every $\beta$ and $\lambda$. Consequently, for bond fields $A^{(1)},\dots,A^{(n)}$ and every $\beta$, the coefficient
\[
W_{n,\beta}(A^{(1)},\dots,A^{(n)};\lambda)\defeq\partial_s\partial_{t_1}\cdots\partial_{t_n}\,[\tau^\beta]\log Z\Big(\sum_it_iA^{(i)}+s\,d\lambda,\tau\Big)\Big|_{t=s=0}
\]
vanishes ($[\tau^\beta]$ denoting the coefficient of $\tau^\beta/\beta!$ in the formal logarithm). $W_{n,\beta}$ is a finite sum of connected graphs with lines $C\defeq H(0)^{-1}$, vertices from the Taylor coefficients of $H(A)$ in $A$ (the vertices \cite[(1.8), (1.10)]{Ba3} with $g=1$ at background $B$) at the bonds of $\supp A^{(i)}$, one vertex linear in $d\lambda$, and $\beta_r$ vertices $F_r$. The identities \cite[(2.25)--(2.28)]{Ba3} are the instances $(n,\beta)=(0,0)$, $(1,0)$, $(1,e_r)$ with $F=-\lambda_k\sum_{x\in\Delta}\eta^d:\!|\phi(x)|^4\!:$ and $B=0$.
\end{lemma}

\begin{proof}
\emph{Finiteness.} $H(A)\ge M^2$ for real $A$, and $\Re H(A)\ge M^2/2$ for $A$ in a complex neighbourhood of the real bond fields; so each $Z_\beta$ is a finite-dimensional Gaussian integral of a polynomial, analytic in $A$ there, and $Z_0>0$ on real $A$. ($\tau$ is a formal parameter only; $e^{\tau F}$ is never integrated.) \emph{Invariance.} Substitute $\phi=T_\lambda^{-1}\psi$: the Jacobian is $1$ (orthogonal), the quadratic form changes by \cref{sup2:lem:S8cov}, and $F^\beta$ is unchanged; so $Z_\beta(A+d\lambda)=Z_\beta(A)$. \emph{Logarithm.} Every coefficient of the formal series $\log Z(A,\tau)$ is a polynomial in the $Z_\beta(A)/Z_0(A)$, hence gauge invariant in $A$, and its $\partial_s$-derivative at $s=0$ vanishes identically in $t$. \emph{Graph form.} By Wick's theorem the $A$-Taylor coefficients at $A=0$ of these coefficients are the connected Gaussian moments of the polynomials $\partial^a_AH(A)|_{A=0}$ and of the $F_r$: finite sums of connected graphs with lines $C$. \emph{Instances.} \cite[(2.25)]{Ba3} is $n=0$; \cite[(2.26)--(2.27)]{Ba3} is $n=1$ at $B=0$ (the two bubbles, the seagull loop and the $\eta$-loop); \cite[(2.28)]{Ba3} adds one $F$-vertex.
\end{proof}

\begin{lemma}[free boundary conditions as the limit of the periodic ones]\label{sup2:lem:S8lim}
On $\eta\Z^d$ let $C_\infty\defeq H(0)^{-1}$ (bounded on $\ell^2$, $H\ge M^2$) with $B$ bounded, let $\lambda$ and the $A^{(i)}$ be finitely supported and the $F_r$ polynomials in the $\phi(x)$, $x\in X_F$ finite. Define $W^\infty_{n,\beta}$ by the finite graph sum of \cref{sup2:lem:S8torus} with $C_\infty$ in place of $C$. Then $W^\infty_{n,\beta}=0$.
\end{lemma}

\begin{proof}
\emph{Finitely many positions.} Every vertex of the graphs sits at a fixed position: the $A^{(i)}$-vertices at bonds of $\supp A^{(i)}$, the $\lambda$-vertex at bonds of $\supp d\lambda$, the $F$-vertices in $X_F$; no position is summed, because $H(A)-H(0)$ vanishes off $\supp A$. So $W^\infty$ is a fixed polynomial in the finitely many numbers $C_\infty(u,u')$, $u,u'\in S$, $S$ the union of the supports and their $\eta$-neighbours. \emph{Embedding in a torus.} Choose $\ell>\operatorname{diam}S+2\eta$ and $B_\ell$ periodic of period $\ell$ agreeing with $B$ on a box of side $\ell$ around $S$ (the periodisation of its restriction; it changes $H$ only at distance $\ge\ell/2-\operatorname{diam}S$ from $S$). The projection $\eta\Z^d\to T_\ell$ is injective on the convex hull of $S$, which contains $\supp\lambda$ (a finitely supported $\lambda$ has its support in the hull of $\supp d\lambda$), so $\lambda$, the $A^{(i)}$ and the $F_r$ lift to $T_\ell$, and \cref{sup2:lem:S8torus} gives $W^{T_\ell}=0$. \emph{Convergence of the kernels.} Let $\tilde C_\infty\defeq(-\Delta_{B_\ell}+M^2)^{-1}$ on $\eta\Z^d$. The torus kernel is the image sum $C_\ell(u,u')=\sum_{n\in\Z^d}\tilde C_\infty(u,u'+n\ell)$: the sum converges absolutely by \cref{sup2:lem:S8CT}, it is $\ell$-periodic, and $(-\Delta_{B_\ell}+M^2)$ applied to it gives $\sum_n\delta_{u'+n\ell}$, so it solves the periodic equation on $T_\ell$, whose solution is unique since $H_\ell\ge M^2>0$. By the second resolvent identity and \cref{sup2:lem:S8CT}, $|\tilde C_\infty-C_\infty|(u,u')\le c'\exp(-\frac\theta2(\ell/2-\operatorname{diam}S))$ on $S\times S$ (one half of each exponential pays the sum over the support of $B_\ell-B$), and the images $n\ne0$ are $\le c\,e^{-\theta(\ell-\operatorname{diam}S)}$ each, summable. \emph{Limit.} Hence $C_\ell(u,u')\to C_\infty(u,u')$ for $u,u'\in S$ as $\ell\to\infty$, and $W^\infty=\lim_\ell W^{T_\ell}=0$, $W$ being a fixed polynomial in these entries.
\end{proof}
